\chapter{The Inflation Flip}

The Inflation Flip is an argumentative behavior pattern in which a specific, bounded grievance is inflated into a sweeping character indictment, and the person who inflated it then positions themselves as the victim of that indictment --- effectively reversing who is accountable to whom.
Whether the person executing it is aware of what they are doing is largely beside the point: the pattern functions the same way regardless of intent, and its effects on the person who raised the original grievance are the same either way.
It is worth understanding in the context of DARVO, a related pattern it is sometimes compared to, though the two are distinct in mechanism and should not be conflated.

\section{DARVO: Brief Background}

DARVO --- Deny, Attack, Reverse Victim and Offender --- is a term coined by psychologist Jennifer Freyd for a defensive pattern in which the person being held accountable denies the behavior, attacks the person raising the grievance, and repositions themselves as the true victim of the exchange.
It is well documented across workplace conflicts, family dynamics, and romantic relationships, and is useful context here primarily because the Inflation Flip shares its terminal state --- role reversal --- while arriving there by a different route.

\section{The Inflation Flip}

The Inflation Flip is a distinct pattern that shares with DARVO only its terminal state: the reversal of victim and offender roles.
It does not require a denial phase, and it does not require an overt attack.
What it requires is inflation and self-victimization --- a different mechanism entirely, and in some ways a more insidious one.

\begin{definition}[The Inflation Flip]
The \emph{Inflation Flip} is a six-step argumentative behavior pattern initiated in response to a specific grievance, proceeding as follows:

\begin{enumerate}
  \item \textbf{Specific grievance.}
  Person A raises a concrete, bounded criticism about a specific action or failure.
  The criticism is narrow and factual.
  \emph{Example: ``You rejected my call twice after I told you it was urgent.''}

  \item \textbf{Overgeneralization.}
  Person B does not engage with the specific criticism.
  Instead, they inflate it into a sweeping indictment of their entire character, history, or identity.
  A single data point becomes an absolute.
  \emph{Example: ``You're saying I'm a terrible friend who is never there for you.''}

  \item \textbf{Strawman construction.}
  The inflated, exaggerated version of the criticism is treated as the actual argument being made.
  Person B is no longer responding to what was said, but to a version they constructed --- one that is far easier to argue against because it is demonstrably false or extreme.

  \item \textbf{Self-victimization.}
  Person B positions themselves as the victim of this strawmanned attack.
  The language shifts to emotional injury: \emph{``It hurts that you'd say that,''} \emph{``I can't believe you think that of me,''} \emph{``You're making me feel like I've never done anything right.''}
  The original grievance has been replaced entirely by Person B's feelings about the grievance.

  \item \textbf{Derailment.}
  Person A is now on the defensive, spending their energy correcting the record and clarifying what they did and did not say.
  The original specific grievance has evaporated.

  \item \textbf{Role reversal.}
  Person B is now the wounded party.
  Person A is the one who must answer for their behavior.
  The positions of accuser and accused have been fully exchanged.
\end{enumerate}
\end{definition}

\begin{remark}[Verb forms and register]
This pattern carries three forms of address, each suited to a different context.

\emph{Inflation flip} (and its verb form, \emph{inflation flipping}) is the formal designation used in definitions, first introductions, and anywhere precision matters.
It names both mechanisms explicitly and leaves no ambiguity.

\emph{Iflip} (verb: \emph{iflipping}, past tense: \emph{iflipped}) is the informal shorthand, suitable once the pattern has been established between speaker and audience.
The prefix \emph{i-} compresses \emph{inflation} and is understood as such by anyone already familiar with the term.
It is short enough to use naturally in conversation and distinct enough that it cannot be confused with anything else.
\begin{itemize}
  \item \emph{``She iflipped me the second I brought it up.''}
  \item \emph{``I didn't realize I was being iflipped until I was already apologizing.''}
  \item \emph{``That's a classic iflip --- he never addressed what you actually said.''}
\end{itemize}

\emph{Flip} (verb: \emph{flipping}, past tense: \emph{flipped}) is the most compressed form and should only be used when the inflation sense is completely unambiguous from context.
It is the most natural in the past tense and present progressive, but risks vagueness if the conversation has drifted or the listener is unfamiliar with the term.

As a rule: use \emph{inflation flip} to introduce, \emph{iflip} to refer, and \emph{flip} only when you are certain it will land.
\end{remark}

Figure~\ref{fig:inflation-flip} illustrates the pipeline.

\begin{figure}[htbp]
  \centering
  \begin{tikzpicture}[
    box/.style={
      draw, rounded corners=4pt,
      minimum width=2.1cm, minimum height=0.8cm,
      align=center, font=\footnotesize
    },
    >=Stealth, thick,
    node distance=0.45cm
  ]
    % Row 1: steps 1–3
    \node[box] (g)  {Specific\\Grievance};
    \node[box, right=of g] (o)  {Over-\\generalization};
    \node[box, right=of o] (s)  {Strawman\\Construction};

    % Row 2: steps 4–6 (right to left so arrows flow downward then left)
    \node[box, below=0.7cm of s] (v)  {Self-\\victimization};
    \node[box, left=of v]        (d)  {Derailment};
    \node[box, left=of d]        (r)  {Role\\Reversal};

    % Arrows row 1
    \draw[->] (g) -- (o);
    \draw[->] (o) -- (s);

    % Turn down
    \draw[->] (s) -- (v);

    % Arrows row 2 (right to left)
    \draw[->] (v) -- (d);
    \draw[->] (d) -- (r);

  \end{tikzpicture}
  \caption{The Inflation Flip pipeline.}
  \label{fig:inflation-flip}
\end{figure}

\section{Relationship to DARVO}

The Inflation Flip and DARVO share a common terminal state: role reversal.
They do not share a mechanism.

DARVO achieves its reversal through denial and direct attack.
The person denies the behavior, goes on the offensive, and positions themselves as the victim of the confrontation itself.
The aggression is often legible --- there is something to push back against.

The Inflation Flip reaches the same endpoint possibly without denial and/or without overt attack.
Person B in the example above does not deny what happened.
They do not accuse Person A of lying or attack their character.
They simply receive the criticism and immediately expand it into something they can suffer from.
The move presents as vulnerability rather than hostility, which is precisely what makes it harder to identify and name in real time.

\begin{remark}
The absence of the Deny and Attack phases is not a mitigating feature --- it is a structural one.
Because the Inflation Flip generates no obvious aggression, the person who raised the original grievance has nothing visible to push back against.
Calling it out requires naming a pattern rather than pointing to an act, which is a harder thing to do in the middle of a conversation.
\end{remark}

\begin{remark}
Not all DARVO involves an Inflation Flip, and not all Inflation Flips involve DARVO.
They are parallel patterns with a shared destination, not nested ones.
Both belong to a broader class of accountability deflection maneuvers in which the terminal state is role reversal.
\end{remark}

\section{Why It Works}

The Inflation Flip exploits several natural human tendencies simultaneously.

Most people do not want to be seen as someone who attacks others unfairly.
When Person B says ``you're making me feel like I've never been a good friend,'' Person A's instinct is to correct that impression --- to reassure, to walk it back.
This is a socially conditioned response to perceived emotional harm.
The Inflation Flip hijacks that instinct and uses it to bury the original issue.

It also exploits the asymmetry between specific and absolute claims.
Specific claims are difficult to deny: ``You didn't answer my call'' is a fact.
Inflated claims, by contrast, are easy to dispute, qualify, and argue against indefinitely --- which is precisely the point.
The moment Person A begins arguing against ``you think I'm always terrible,'' they have left the terrain of fact and entered the terrain of characterization.
That terrain favors Person B, who constructed it.

% \begin{proposition}
% Let $G_s$ denote a specific, bounded grievance and let $G_g = I(G_s)$ denote its inflated generalization under the Inflation Flip.
% The following asymmetry holds:
% \begin{itemize}
%   \item $G_s$ is difficult to deny because it is grounded in a particular action at a particular time.
%   \item $G_g$ is easy to dispute because it is an absolute claim about character or pattern, which admits counterexamples and ambiguity.
% \end{itemize}
% The substitution of $G_g$ for $G_s$ as the subject of the exchange is therefore strategically advantageous to Person B regardless of whether the inflation is conscious.
% \end{proposition}

\section{Examples}

\paragraph{Friendship conflict.}
\begin{quote}
\textbf{A:} ``You didn't show up when I needed you and I told you it was urgent.''\\
\textbf{B:} ``So you're saying I'm a terrible friend who has never been there for you. I can't believe after everything I've done you'd make me feel like I've never done anything right.''
\end{quote}
Person A is now clarifying and apologizing.
The missed call is no longer the subject.

\paragraph{Workplace.}
\begin{quote}
\textbf{A:} ``You didn't credit me in that presentation.''\\
\textbf{B:} ``Wow, so you think I always steal credit and I've never once acknowledged your work. That's really hurtful to hear.''
\end{quote}
Person A is now defending whether they really think that, not addressing the presentation.

\paragraph{Romantic relationship.}
\begin{quote}
\textbf{A:} ``You were on your phone the entire dinner.''\\
\textbf{B:} ``You're saying I'm emotionally unavailable and I never pay attention to you. I give so much to this relationship and this is what you think of me.''
\end{quote}
Person A is now managing Person B's feelings about a characterization Person A never made.

\section{Identifying It in Real Time}

\begin{remark}[Signal: absolute language]
The clearest signal is the sudden appearance of absolute quantifiers: \emph{always}, \emph{never}, \emph{every time}, \emph{not once}.
These quantifiers almost never appear in the original grievance but emerge immediately in the inflated response.
When they appear, it is a reliable indicator that the specific has been inflated into the general and the Inflation Flip is underway.
\end{remark}

\begin{remark}[Signal: role inversion]
The second signal is the sudden shift in who is expressing hurt.
If you raised a grievance and within one or two exchanges you find yourself apologizing and explaining, ask: \emph{what am I apologizing for?}
If the answer is something you did not actually say or claim, the flip has already occurred.
\end{remark}

